\scanpage{14}
$k$ anneau (com. unit.)

\uline{Déf 1} Un $k$-module $\frac12$ simplicial \inserted{$L_\bullet$} est dit $\partial$-pseudo-cohérent (resp. strictement $\partial$-pseudocohérent, resp. $\partial$-parfait, resp strictement $\partial$parfait) si $L_\bullet^!$ (le complexe de chaînes associé par Dold-Puppe) est pseudo-cohérent (resp. str. ps. cohérent, resp parfait, resp. strictt parfait). Donc \unsure{de-par|de par} déf a) \inserted{ci-dessous}

\uline{Prop 1} a) $L_\bullet$ est str. $\deleted{\unsure{\phi}}\partial$-pseudocohérent (resp. str. $\partial$-parfait) si et seulmt si les $(L_n^!)_{\unsure{n}}$ sont projectifs de t. f. (resp. et nuls sauf un nb fini), \deleted{\illegible{2}}

b) $L_\bullet$ est str. $\partial$-ps. coh. \unsure{sss|ssi} les $L_n$ sont projectifs de t. f.

c) $L_\bullet$ est str. $\partial$parfait sss $L_\bullet$ est extension successive \inserted{finie} de Modules de la forme $\Lambda^i\Psi_\bullet\otimes_{\mathbb Z}M_i$ \inserted{$(i\geq0)$}, où les $M_i$ sont $k$-modules pr. de t.f. (l'ordre des facteurs dans une filtration croissante \unsure{peut être pris|peuvent être|pouvant} d'ailleurs être choisi par ordre des $i$ croissants), ou encore \deleted{\unsure{sss|ssi}} \inserted{(ce qui suit se \unsure{comprend} cas spécial : commencer par s'y placer)}
% slanted
\margin{sss les $L_n$ sont proj. de t. f. et $f(n)=\mathrm{rg}_k L_n$ est une fonction polynômiale pour $n$ grand (on montrera \unsure{ici} une telle fonction)}

d) Si $k$ est noeth., $L_\bullet$ est $\partial$-pseudo-cohérent (sss les $L_n^!$ sont de type fini sur $k$, ou encore) sss les $L_n^{\unsure{\flat|\natural}}$ le sont.

e) Si $k$ est noeth. régulier, $L_\bullet$ est $\partial$-parfait sss les $L_n$ sont de t. f. et \deleted{\unsure{la} $\unsure{\ldots}$ \unsure{fn.}} $\forall\,\unsure{k|K}$ corps résiduel, la fonction
% archive sheet number: 13
